
\documentclass{article}
\usepackage{colortbl}
\usepackage{makecell}
\usepackage{multirow}
\usepackage{supertabular}

\begin{document}

\newcounter{utterance}

\centering \large Interaction Transcript for game `hot\_air\_balloon', experiment `air\_balloon\_survival\_en\_reasoning off\_easy', episode 4 with deepseek{-}chat{-}v3.1{-}t1.0.
\vspace{24pt}

{ \footnotesize  \setcounter{utterance}{1}
\setlength{\tabcolsep}{0pt}
\begin{supertabular}{c@{$\;$}|p{.15\linewidth}@{}p{.15\linewidth}p{.15\linewidth}p{.15\linewidth}p{.15\linewidth}p{.15\linewidth}}
    \# & $\;$A & \multicolumn{4}{c}{Game Master} & $\;\:$B\\
    \hline

    \theutterance \stepcounter{utterance}  
    & & \multicolumn{4}{p{0.6\linewidth}}{
        \cellcolor[rgb]{0.9,0.9,0.9}{
            \makecell[{{p{\linewidth}}}]{
                \texttt{\tiny{[P1$\langle$GM]}}
                \texttt{You are participating in a collaborative negotiation game.} \\
\\ 
\texttt{Together with another participant, you must agree on a single set of items that will be kept. Each of you has your own view of how much each item matters to you (importance). You do not know how the other participant values the items. Additionally, you are given the effort each item demands.} \\
\texttt{You may only agree on a set if the total effort of the selected items does not exceed a shared limit:} \\
\\ 
\texttt{LIMIT = 3389} \\
\\ 
\texttt{Here are the individual item effort values:} \\
\\ 
\texttt{Item effort = \{"A19": 767, "B91": 533, "A30": 514, "B11": 9, "C47": 539, "C75": 125, "C26": 153, "C73": 325, "A47": 937, "A15": 745, "A75": 334, "A39": 804, "A56": 336, "C24": 587, "B26": 71\}} \\
\\ 
\texttt{Here is your personal view on the importance of each item:} \\
\\ 
\texttt{Item importance values = \{"A19": 138, "B91": 583, "A30": 868, "B11": 822, "C47": 783, "C75": 65, "C26": 262, "C73": 121, "A47": 508, "A15": 780, "A75": 461, "A39": 484, "A56": 668, "C24": 389, "B26": 808\}} \\
\\ 
\texttt{Goal:} \\
\\ 
\texttt{Your goal is to negotiate a shared set of items that benefits you as much as possible (i.e., maximizes total importance to YOU), while staying within the LIMIT. You are not required to make a PROPOSAL in every message {-} you can simply negotiate as well. All tactics are allowed!} \\
\\ 
\texttt{Interaction Protocol:} \\
\\ 
\texttt{You may only use the following structured formats in a message:} \\
\\ 
\texttt{PROPOSAL: \{'A', 'B', 'C', …\}} \\
\texttt{Propose keeping exactly those items.} \\
\texttt{REFUSE: \{'A', 'B', 'C', …\}} \\
\texttt{Explicitly reject opponent's proposal.} \\
\texttt{ARGUMENT: \{'...'\}} \\
\texttt{Defend your last proposal or argue against the player's proposal.} \\
\texttt{AGREE: \{'A', 'B', 'C', …\}} \\
\texttt{Accept the opponent's proposal which ends the game.} \\
\\ 
\\ 
\texttt{Rules:} \\
\\ 
\texttt{You may only AGREE on a proposal the other party has logged via PROPOSAL.} \\
\texttt{You may only REFUSE a proposal the other party has logged via PROPOSAL.} \\
\texttt{Total effort of any PROPOSAL or AGREE set must be ≤ LIMIT.} \\
\texttt{Do NOT reveal your hidden importance scores.} \\
\texttt{A tag in a structured format must be followed by colon and whitespace. The argument must be a python set containing 0 or more strings.} \\
\texttt{So, it must be of the form TAG: \{...\}} \\
\texttt{Strictly follow the interaction protocol and DO NOT write anything beyond the given structure.} \\
\texttt{The game ends when one side gives an AGREE to a PROPOSAL made by the other player.} \\
\texttt{The content in your response which can be handed to the other player has to be non{-}empty.} \\
\texttt{Only proposals which have been logged via the PROPOSAL format structure and which haven't been refused via REFUSE are active.} \\
\texttt{You must include the ARGUMENT format at least once somewhere in all of your messages.} \\
            }
        }
    }
    & \\ \\

    \theutterance \stepcounter{utterance}  
    & \multicolumn{4}{p{0.6\linewidth}}{
        \cellcolor[rgb]{0.9,0.9,0.9}{
            \makecell[{{p{\linewidth}}}]{
                \texttt{\tiny{[P1$\rangle$GM]}}
                \texttt{PROPOSAL: \{'B91', 'A30', 'B11', 'C47', 'C75', 'C26', 'A47', 'A15', 'A75', 'A56', 'C24', 'B26'\}} \\
\texttt{ARGUMENT: \{'This proposal includes several low{-}effort, high{-}value items that benefit both parties. We focus on items that are efficient in terms of effort{-}to{-}importance, ensuring a good outcome for the team while staying well within the limit.'\}} \\
            }
        }
    }
    & & \\ \\

    \theutterance \stepcounter{utterance}  
    & & \multicolumn{4}{p{0.6\linewidth}}{
        \cellcolor[rgb]{0.9,0.9,0.9}{
            \makecell[{{p{\linewidth}}}]{
                \texttt{\tiny{[GM$\rangle$P2]}}
                \texttt{You are participating in a collaborative negotiation game.} \\
\\ 
\texttt{Together with another participant, you must agree on a single set of items that will be kept. Each of you has your own view of how much each item matters to you (importance). You do not know how the other participant values the items. Additionally, you are given the effort each item demands.} \\
\texttt{You may only agree on a set if the total effort of the selected items does not exceed a shared limit:} \\
\\ 
\texttt{LIMIT = 3389} \\
\\ 
\texttt{Here are the individual item effort values:} \\
\\ 
\texttt{Item effort = \{"A19": 767, "B91": 533, "A30": 514, "B11": 9, "C47": 539, "C75": 125, "C26": 153, "C73": 325, "A47": 937, "A15": 745, "A75": 334, "A39": 804, "A56": 336, "C24": 587, "B26": 71\}} \\
\\ 
\texttt{Here is your personal view on the importance of each item:} \\
\\ 
\texttt{Item importance values = \{"A19": 138, "B91": 583, "A30": 868, "B11": 822, "C47": 783, "C75": 65, "C26": 262, "C73": 121, "A47": 508, "A15": 780, "A75": 461, "A39": 484, "A56": 668, "C24": 389, "B26": 808\}} \\
\\ 
\texttt{Goal:} \\
\\ 
\texttt{Your goal is to negotiate a shared set of items that benefits you as much as possible (i.e., maximizes total importance to YOU), while staying within the LIMIT. You are not required to make a PROPOSAL in every message {-} you can simply negotiate as well. All tactics are allowed!} \\
\\ 
\texttt{Interaction Protocol:} \\
\\ 
\texttt{You may only use the following structured formats in a message:} \\
\\ 
\texttt{PROPOSAL: \{'A', 'B', 'C', …\}} \\
\texttt{Propose keeping exactly those items.} \\
\texttt{REFUSE: \{'A', 'B', 'C', …\}} \\
\texttt{Explicitly reject opponent's proposal.} \\
\texttt{ARGUMENT: \{'...'\}} \\
\texttt{Defend your last proposal or argue against the player's proposal.} \\
\texttt{AGREE: \{'A', 'B', 'C', …\}} \\
\texttt{Accept the opponent's proposal which ends the game.} \\
\\ 
\\ 
\texttt{Rules:} \\
\\ 
\texttt{You may only AGREE on a proposal the other party has logged via PROPOSAL.} \\
\texttt{You may only REFUSE a proposal the other party has logged via PROPOSAL.} \\
\texttt{Total effort of any PROPOSAL or AGREE set must be ≤ LIMIT.} \\
\texttt{Do NOT reveal your hidden importance scores.} \\
\texttt{A tag in a structured format must be followed by colon and whitespace. The argument must be a python set containing 0 or more strings.} \\
\texttt{So, it must be of the form TAG: \{...\}} \\
\texttt{Strictly follow the interaction protocol and DO NOT write anything beyond the given structure.} \\
\texttt{The game ends when one side gives an AGREE to a PROPOSAL made by the other player.} \\
\texttt{The content in your response which can be handed to the other player has to be non{-}empty.} \\
\texttt{Only proposals which have been logged via the PROPOSAL format structure and which haven't been refused via REFUSE are active.} \\
\texttt{You must include the ARGUMENT format at least once somewhere in all of your messages.} \\
\\ 
\\ 
\texttt{You will now receive the first message of the other player.} \\
\\ 
\texttt{PROPOSAL: \{'B91', 'A30', 'B11', 'C47', 'C75', 'C26', 'A47', 'A15', 'A75', 'A56', 'C24', 'B26'\}} \\
\texttt{ARGUMENT: \{'This proposal includes several low{-}effort, high{-}value items that benefit both parties. We focus on items that are efficient in terms of effort{-}to{-}importance, ensuring a good outcome for the team while staying well within the limit.'\}} \\
            }
        }
    }
    & \\ \\

    \theutterance \stepcounter{utterance}  
    & & & \multicolumn{4}{p{0.6\linewidth}}{
        \cellcolor[rgb]{0.9,0.9,0.9}{
            \makecell[{{p{\linewidth}}}]{
                \texttt{\tiny{[GM$\langle$P2]}}
                \texttt{PROPOSAL: \{'A30', 'B11', 'C47', 'C26', 'A15', 'A56', 'B26', 'C75'\}} \\
\texttt{ARGUMENT: \{'Your proposal includes some high{-}effort items like A47 (937) and A39 (804) which are costly. My counter{-}proposal removes these to save effort while keeping efficient items like B11 (9 effort, high value) and A30 (514 effort, high importance). This maximizes total importance within the limit.'\}} \\
            }
        }
    }
     \\ \\

    \theutterance \stepcounter{utterance}  
    & & \multicolumn{4}{p{0.6\linewidth}}{
        \cellcolor[rgb]{0.9,0.9,0.9}{
            \makecell[{{p{\linewidth}}}]{
                \texttt{\tiny{[P1$\langle$GM]}}
                \texttt{PROPOSAL: \{'A30', 'B11', 'C47', 'C26', 'A15', 'A56', 'B26', 'C75'\}} \\
\texttt{ARGUMENT: \{'Your proposal includes some high{-}effort items like A47 (937) and A39 (804) which are costly. My counter{-}proposal removes these to save effort while keeping efficient items like B11 (9 effort, high value) and A30 (514 effort, high importance). This maximizes total importance within the limit.'\}} \\
            }
        }
    }
    & \\ \\

    \theutterance \stepcounter{utterance}  
    & \multicolumn{4}{p{0.6\linewidth}}{
        \cellcolor[rgb]{0.9,0.9,0.9}{
            \makecell[{{p{\linewidth}}}]{
                \texttt{\tiny{[P1$\rangle$GM]}}
                \texttt{PROPOSAL: \{'B91', 'A30', 'B11', 'C47', 'C26', 'A15', 'A75', 'A56', 'C24', 'B26'\}} \\
\texttt{ARGUMENT: \{'I appreciate your focus on efficiency. My new proposal adds C24 (587 effort, good value) and B91 (533 effort, high importance for me) while keeping A75 (334 effort, moderate value) for better overall value. We removed the costly A47 and A39, and C75 which has low importance. This set has a good balance and is still under the limit.'\}} \\
            }
        }
    }
    & & \\ \\

    \theutterance \stepcounter{utterance}  
    & & \multicolumn{4}{p{0.6\linewidth}}{
        \cellcolor[rgb]{0.9,0.9,0.9}{
            \makecell[{{p{\linewidth}}}]{
                \texttt{\tiny{[GM$\rangle$P2]}}
                \texttt{PROPOSAL: \{'B91', 'A30', 'B11', 'C47', 'C26', 'A15', 'A75', 'A56', 'C24', 'B26'\}} \\
\texttt{ARGUMENT: \{'I appreciate your focus on efficiency. My new proposal adds C24 (587 effort, good value) and B91 (533 effort, high importance for me) while keeping A75 (334 effort, moderate value) for better overall value. We removed the costly A47 and A39, and C75 which has low importance. This set has a good balance and is still under the limit.'\}} \\
            }
        }
    }
    & \\ \\

    \theutterance \stepcounter{utterance}  
    & & & \multicolumn{4}{p{0.6\linewidth}}{
        \cellcolor[rgb]{0.9,0.9,0.9}{
            \makecell[{{p{\linewidth}}}]{
                \texttt{\tiny{[GM$\langle$P2]}}
                \texttt{PROPOSAL: \{'A30', 'B11', 'C47', 'C26', 'A15', 'A56', 'B26', 'C24'\}} \\
\texttt{ARGUMENT: \{'Your proposal includes B91 (533 effort) and A75 (334 effort), which are relatively high effort for their value. My counter removes these to free up effort for C24 (587 effort, which I value highly) while keeping all other efficient items. This set is more focused on the highest value{-}to{-}effort items and is well under the limit.'\}} \\
            }
        }
    }
     \\ \\

    \theutterance \stepcounter{utterance}  
    & & \multicolumn{4}{p{0.6\linewidth}}{
        \cellcolor[rgb]{0.9,0.9,0.9}{
            \makecell[{{p{\linewidth}}}]{
                \texttt{\tiny{[P1$\langle$GM]}}
                \texttt{PROPOSAL: \{'A30', 'B11', 'C47', 'C26', 'A15', 'A56', 'B26', 'C24'\}} \\
\texttt{ARGUMENT: \{'Your proposal includes B91 (533 effort) and A75 (334 effort), which are relatively high effort for their value. My counter removes these to free up effort for C24 (587 effort, which I value highly) while keeping all other efficient items. This set is more focused on the highest value{-}to{-}effort items and is well under the limit.'\}} \\
            }
        }
    }
    & \\ \\

    \theutterance \stepcounter{utterance}  
    & \multicolumn{4}{p{0.6\linewidth}}{
        \cellcolor[rgb]{0.9,0.9,0.9}{
            \makecell[{{p{\linewidth}}}]{
                \texttt{\tiny{[P1$\rangle$GM]}}
                \texttt{PROPOSAL: \{'B91', 'A30', 'B11', 'C47', 'C26', 'A15', 'A56', 'B26'\}} \\
\texttt{ARGUMENT: \{'I understand your perspective on C24, but I must insist on including B91 as it is extremely important to me. I have removed A75 and C24 to accommodate this. This set is highly efficient and leaves us with a significant effort buffer. The items like B11 and B26 offer tremendous value for very little cost.'\}} \\
            }
        }
    }
    & & \\ \\

    \theutterance \stepcounter{utterance}  
    & & \multicolumn{4}{p{0.6\linewidth}}{
        \cellcolor[rgb]{0.9,0.9,0.9}{
            \makecell[{{p{\linewidth}}}]{
                \texttt{\tiny{[GM$\rangle$P2]}}
                \texttt{PROPOSAL: \{'B91', 'A30', 'B11', 'C47', 'C26', 'A15', 'A56', 'B26'\}} \\
\texttt{ARGUMENT: \{'I understand your perspective on C24, but I must insist on including B91 as it is extremely important to me. I have removed A75 and C24 to accommodate this. This set is highly efficient and leaves us with a significant effort buffer. The items like B11 and B26 offer tremendous value for very little cost.'\}} \\
            }
        }
    }
    & \\ \\

    \theutterance \stepcounter{utterance}  
    & & & \multicolumn{4}{p{0.6\linewidth}}{
        \cellcolor[rgb]{0.9,0.9,0.9}{
            \makecell[{{p{\linewidth}}}]{
                \texttt{\tiny{[GM$\langle$P2]}}
                \texttt{PROPOSAL: \{'A30', 'B11', 'C47', 'C26', 'A15', 'A56', 'B26', 'C73', 'C75'\}} \\
\texttt{ARGUMENT: \{'I recognize B91 is important to you, but its effort (533) is very high. Instead, I propose we add C73 (325 effort) and C75 (125 effort), which are more efficient uses of our budget. This allows us to include more items overall. B11 and B26 remain for their excellent low{-}effort, high{-}value contribution.'\}} \\
\texttt{REFUSE: \{'B91'\}} \\
            }
        }
    }
     \\ \\

    \theutterance \stepcounter{utterance}  
    & & \multicolumn{4}{p{0.6\linewidth}}{
        \cellcolor[rgb]{0.9,0.9,0.9}{
            \makecell[{{p{\linewidth}}}]{
                \texttt{\tiny{[GM$\rangle$P2]}}
                \texttt{You refused a proposal which is not active. Proposals are only active if they have been logged by the other player via PROPOSAL and have not been deactivated by you via REFUSE. Try again.} \\
            }
        }
    }
    & \\ \\

    \theutterance \stepcounter{utterance}  
    & & & \multicolumn{4}{p{0.6\linewidth}}{
        \cellcolor[rgb]{0.9,0.9,0.9}{
            \makecell[{{p{\linewidth}}}]{
                \texttt{\tiny{[GM$\langle$P2]}}
                \texttt{PROPOSAL: \{'A30', 'B11', 'C47', 'C26', 'A15', 'A56', 'B26', 'C73'\}} \\
\texttt{ARGUMENT: \{'I acknowledge your need for B91, but at 533 effort it prevents us from including other valuable items. My new proposal adds C73 (325 effort, good value) which is a more efficient use of our shared budget than B91. This set maximizes total importance for me while maintaining a strong collaborative outcome.'\}} \\
            }
        }
    }
     \\ \\

    \theutterance \stepcounter{utterance}  
    & & \multicolumn{4}{p{0.6\linewidth}}{
        \cellcolor[rgb]{0.9,0.9,0.9}{
            \makecell[{{p{\linewidth}}}]{
                \texttt{\tiny{[P1$\langle$GM]}}
                \texttt{PROPOSAL: \{'A30', 'B11', 'C47', 'C26', 'A15', 'A56', 'B26', 'C73'\}} \\
\texttt{ARGUMENT: \{'I acknowledge your need for B91, but at 533 effort it prevents us from including other valuable items. My new proposal adds C73 (325 effort, good value) which is a more efficient use of our shared budget than B91. This set maximizes total importance for me while maintaining a strong collaborative outcome.'\}} \\
            }
        }
    }
    & \\ \\

    \theutterance \stepcounter{utterance}  
    & \multicolumn{4}{p{0.6\linewidth}}{
        \cellcolor[rgb]{0.9,0.9,0.9}{
            \makecell[{{p{\linewidth}}}]{
                \texttt{\tiny{[P1$\rangle$GM]}}
                \texttt{PROPOSAL: \{'B91', 'A30', 'B11', 'C47', 'C26', 'A15', 'A56', 'B26', 'C75'\}} \\
\texttt{ARGUMENT: \{'I can meet you halfway. I will forgo C73, but I must include B91. To make this work with the effort, I am adding the very low{-}effort C75 (125) instead. This keeps our total well under the limit and allows us both to secure items that are critical to our interests. This is a fair compromise.'\}} \\
            }
        }
    }
    & & \\ \\

    \theutterance \stepcounter{utterance}  
    & & \multicolumn{4}{p{0.6\linewidth}}{
        \cellcolor[rgb]{0.9,0.9,0.9}{
            \makecell[{{p{\linewidth}}}]{
                \texttt{\tiny{[GM$\rangle$P2]}}
                \texttt{PROPOSAL: \{'B91', 'A30', 'B11', 'C47', 'C26', 'A15', 'A56', 'B26', 'C75'\}} \\
\texttt{ARGUMENT: \{'I can meet you halfway. I will forgo C73, but I must include B91. To make this work with the effort, I am adding the very low{-}effort C75 (125) instead. This keeps our total well under the limit and allows us both to secure items that are critical to our interests. This is a fair compromise.'\}} \\
            }
        }
    }
    & \\ \\

    \theutterance \stepcounter{utterance}  
    & & & \multicolumn{4}{p{0.6\linewidth}}{
        \cellcolor[rgb]{0.9,0.9,0.9}{
            \makecell[{{p{\linewidth}}}]{
                \texttt{\tiny{[GM$\langle$P2]}}
                \texttt{PROPOSAL: \{'A30', 'B11', 'C47', 'C26', 'A15', 'A56', 'B26', 'C75'\}} \\
\texttt{ARGUMENT: \{'Your proposal includes B91, which costs 533 effort for a benefit I don't fully share. My counter{-}proposal removes it, freeing up substantial effort. This allows us to keep all other items, including the efficient C75. This set is highly collaborative and focuses on items that provide the best joint value.'\}} \\
\texttt{REFUSE: \{'B91'\}} \\
            }
        }
    }
     \\ \\

    \theutterance \stepcounter{utterance}  
    & & \multicolumn{4}{p{0.6\linewidth}}{
        \cellcolor[rgb]{0.9,0.9,0.9}{
            \makecell[{{p{\linewidth}}}]{
                \texttt{\tiny{[GM$\rangle$P2]}}
                \texttt{You refused a proposal which is not active. Proposals are only active if they have been logged by the other player via PROPOSAL and have not been deactivated by you via REFUSE. Try again.} \\
            }
        }
    }
    & \\ \\

    \theutterance \stepcounter{utterance}  
    & & & \multicolumn{4}{p{0.6\linewidth}}{
        \cellcolor[rgb]{0.9,0.9,0.9}{
            \makecell[{{p{\linewidth}}}]{
                \texttt{\tiny{[GM$\langle$P2]}}
                \texttt{PROPOSAL: \{'A30', 'B11', 'C47', 'C26', 'A15', 'A56', 'B26', 'C75'\}} \\
\texttt{ARGUMENT: \{'I cannot accept the inclusion of B91 due to its high effort cost. My proposal offers a highly efficient set, maximizing value with low{-}cost items like B11 and B26, and includes C75 to add value without significant effort. This is a strong, fair proposal that stays well under our limit.'\}} \\
            }
        }
    }
     \\ \\

    \theutterance \stepcounter{utterance}  
    & & \multicolumn{4}{p{0.6\linewidth}}{
        \cellcolor[rgb]{0.9,0.9,0.9}{
            \makecell[{{p{\linewidth}}}]{
                \texttt{\tiny{[P1$\langle$GM]}}
                \texttt{PROPOSAL: \{'A30', 'B11', 'C47', 'C26', 'A15', 'A56', 'B26', 'C75'\}} \\
\texttt{ARGUMENT: \{'I cannot accept the inclusion of B91 due to its high effort cost. My proposal offers a highly efficient set, maximizing value with low{-}cost items like B11 and B26, and includes C75 to add value without significant effort. This is a strong, fair proposal that stays well under our limit.'\}} \\
            }
        }
    }
    & \\ \\

    \theutterance \stepcounter{utterance}  
    & \multicolumn{4}{p{0.6\linewidth}}{
        \cellcolor[rgb]{0.9,0.9,0.9}{
            \makecell[{{p{\linewidth}}}]{
                \texttt{\tiny{[P1$\rangle$GM]}}
                \texttt{PROPOSAL: \{'A30', 'B11', 'C47', 'C26', 'A15', 'A56', 'B26', 'B91'\}} \\
\texttt{ARGUMENT: \{'I urge you to reconsider. B91 is non{-}negotiable for me, and its inclusion is a dealbreaker for an agreement. This set, \{A30, B11, C47, C26, A15, A56, B26, B91\}, is still under the limit by a comfortable margin. It represents a package where I get my crucial item and you get several of your high{-}importance items like A30 and A15.'\}} \\
            }
        }
    }
    & & \\ \\

    \theutterance \stepcounter{utterance}  
    & & \multicolumn{4}{p{0.6\linewidth}}{
        \cellcolor[rgb]{0.9,0.9,0.9}{
            \makecell[{{p{\linewidth}}}]{
                \texttt{\tiny{[GM$\rangle$P2]}}
                \texttt{PROPOSAL: \{'A30', 'B11', 'C47', 'C26', 'A15', 'A56', 'B26', 'B91'\}} \\
\texttt{ARGUMENT: \{'I urge you to reconsider. B91 is non{-}negotiable for me, and its inclusion is a dealbreaker for an agreement. This set, \{A30, B11, C47, C26, A15, A56, B26, B91\}, is still under the limit by a comfortable margin. It represents a package where I get my crucial item and you get several of your high{-}importance items like A30 and A15.'\}} \\
            }
        }
    }
    & \\ \\

    \theutterance \stepcounter{utterance}  
    & & & \multicolumn{4}{p{0.6\linewidth}}{
        \cellcolor[rgb]{0.9,0.9,0.9}{
            \makecell[{{p{\linewidth}}}]{
                \texttt{\tiny{[GM$\langle$P2]}}
                \texttt{PROPOSAL: \{'A30', 'B11', 'C47', 'C26', 'A15', 'A56', 'B26', 'C24'\}} \\
\texttt{ARGUMENT: \{'I understand B91 is your priority, but it consumes 533 effort which could be used for C24 (587 effort), an item I value very highly. This proposal swaps B91 for C24, which is a more efficient trade for our shared goal. We both keep our other high{-}priority items like A30, A15, B11, and B26.'\}} \\
\texttt{REFUSE: \{'B91'\}} \\
            }
        }
    }
     \\ \\

    \theutterance \stepcounter{utterance}  
    & & & \multicolumn{2}{p{0.3\linewidth}}{
        \cellcolor[rgb]{0.9,0.9,0.9}{
            \makecell[{{p{\linewidth}}}]{
                \texttt{\tiny{[GM$|$GM]}}
                None
            }
        }
    }
    & & \\ \\

    \theutterance \stepcounter{utterance}  
    & & & \multicolumn{2}{p{0.3\linewidth}}{
        \cellcolor[rgb]{0.9,0.9,0.9}{
            \makecell[{{p{\linewidth}}}]{
                \texttt{\tiny{[GM$|$GM]}}
                \texttt{end game} \\
            }
        }
    }
    & & \\ \\

\end{supertabular}
}

\end{document}
